\subsection{Equity and the Digital Divide}
\label{sec:8.7.9}
AI development can narrow global inequality or widen it, and which one happens is not automatic. Diagnostic tools that work as well in an under-resourced clinic as in a well-funded hospital, agricultural monitoring that helps a smallholder farmer the way it helps an industrial one, and translation tools that make information equally available regardless of the language a person reads in are all available in principle, and none of them arrives by default. A model trained overwhelmingly on data from wealthy, English-speaking, urban contexts underperforms everywhere else, and a government or company with no reason to fix that gap usually does not.

\begin{esbox}
\boxtitle{The scale of the connectivity gap, 2025}
The International Telecommunication Union's 2025 assessment put global internet use at roughly six billion people, about three-quarters of the world's population, with 2.2 billion people offline, concentrated overwhelmingly in low- and middle-income countries. The gap was starker in quality than in bare access: 94 percent of people in high-income countries used the internet against 23 percent in low-income ones, and 5G coverage reached 84 percent of high-income populations against 4 percent of low-income ones \autocite{itu2025facts}.

\end{esbox}
An AI system trained and deployed for the three-quarters of the world already online is not, without deliberate correction, a system built for the rest. The constraints are structural. Reliable data, computing infrastructure, and the specialized personnel needed to build and maintain AI systems are unevenly distributed, and a country short on all three cannot simply adopt a system built elsewhere and expect it to transfer cleanly. Weak regulatory capacity compounds the problem: a country able to build strong AI safeguards is generally also a country able to enforce them, and the states least equipped to build AI are correspondingly the states least equipped to govern whatever gets deployed within their borders, regardless of who built it. The export controls discussed above cut both ways here: a control written to stop a surveillance system from reaching an authoritarian buyer can just as easily gatekeep the same underlying technology away from a legitimate one, and there is no way to write the rule so it only ever blocks the harmful case.

The multilateral response so far is a commitment and not a mechanism. The Global Digital Compact, adopted by all 193 UN member states at the 2024 Summit of the Future, is the broadest such commitment to closing the digital divide alongside AI governance specifically, and it created two follow-on bodies to carry it: an independent scientific panel and a standing intergovernmental dialogue \autocite{un2024globaldigitalcompact}. What it did not create is an obligation, which puts it in the same category as the voluntary instruments earlier in this section and subject to the same discount.

The remedies proposed at this level have been the same five for years — capacity building that trains local expertise instead of importing it, financing that does not arrive on private investment's terms, intellectual-property flexibility, technology transfer structured as collaboration, and research institutions based in the countries the systems are meant to serve. Every item is correct and none of them is contested. They have been named, in those words, in successive declarations across the whole period the box above measures, and the box is what happened anyway. That is this section's argument reproduced at development scale: the commitment was never the scarce thing. A government or company treating AI primarily as strategic advantage has no reason to fund any of the five, and a declaration is not a reason.

\runin{What This Geopolitics Means for the Rest of the Book}
Every argument this book makes about building an AI system that resists authoritarian capture assumes a government somewhere is deciding whether to let that system be built, sold, or deployed as designed. Sections~\ref{sec:8.7.6} through \ref{sec:8.7.9} are an account of who those governments are, what tools they actually have, and how unevenly those tools are distributed: a handful of companies and the states that host them control the compute the most capable systems require; a binding international treaty exists but binds almost no one yet; voluntary standards multiply faster than any one of them gains real authority; and the states with the weakest capacity to build AI are also the states with the weakest capacity to govern it once someone else's AI arrives at their border, digital or physical. None of this is a reason to abandon the design-level arguments the rest of this book makes — a system built resistant to authoritarian misuse is more resistant regardless of which government eventually asks it to misbehave — but it is a reason not to expect law, international or domestic, to do work that only a system's own design can do.
